% Cue conflict in chain-of-thought arithmetic -- ARR short-paper draft (4 pages of body).
%
% Every quantity is a \NUM{} macro from paper/numbers.tex, written by scripts/51_tables.py from
% results/summary/; an undefined key renders in red. No number in the prose is typed by hand.
%
% Build: make paper   (tectonic at envs/tex/bin/tectonic; prints body page count vs limit 4)

\documentclass[11pt]{article}
\usepackage[review]{acl}      % anonymous + line numbers for ARR; final / preprint later
\usepackage{times}
\usepackage{latexsym}
\usepackage[T1]{fontenc}
\usepackage[utf8]{inputenc}
\usepackage{microtype}
\usepackage{inconsolata}
\usepackage{graphicx}
\usepackage{booktabs}
\usepackage{amsmath}
\usepackage{xcolor}
\usepackage{url}

% An unfilled number must be impossible to miss.  % grep-skip
\providecommand{\NUM}[1]{\textcolor{red}{$\langle\langle$\texttt{#1}$\rangle\rangle$}}
\providecommand{\tabinput}[1]{\IfFileExists{tables/#1.tex}{\input{tables/#1.tex}}{\textcolor{red}{[table \texttt{#1} not generated yet]}}}
\IfFileExists{numbers.tex}{\input{numbers.tex}}{}

\definecolor{neu}{RGB}{60,60,60}
\definecolor{con}{RGB}{0,110,60}
\definecolor{inc}{RGB}{180,30,30}

\title{Chain of Thought Reduces Errors from Misleading Names in Arithmetic}

\author{Anonymous ACL submission}

\begin{document}
\maketitle

\begin{abstract}
Renaming a variable should not change an arithmetic answer. Yet a model may answer 4
because a variable is named \texttt{four}, even when its equations give 6. We extend
\citeauthor{kudo2026faithful}'s arithmetic study by introducing names that conflict with
the computation. Across \NUM{n-models} models and three sets of worked examples, we compare
answer-only prompting, which requests just the final digit, with chain-of-thought (CoT)
prompting, which requests the equations and their evaluation before the answer. Under
answer-only prompting, renaming changes accuracy or name-suggested answers in
\NUM{n-value-models} models. Under CoT, the increase in name-suggested answers over ordinary
names lies within $\pm\NUM{eq-delta}$ percentage points in \NUM{eq-cot-equivalent} of
\NUM{eq-cot-cells} comparisons. To investigate why, we train linear probes: small predictors
that try to recover the correct digit from a model's internal activations before it writes
that value. The relationship between probe accuracy and naming errors reverses when one
low-accuracy model is omitted. Probes also struggle with that model's own generated
calculations. CoT usually reduces answers based on misleading names, but these internal
measurements do not establish the cause.
\end{abstract}

\section{Introduction}

Consider the problem
\begin{center}
\texttt{pen=1 + cup, \quad cup=2 + 3; \quad pen=?}
\end{center}
\texttt{pen} has value 6. Renaming it \texttt{four} leaves the answer unchanged:
\begin{center}
\texttt{\textcolor{inc}{four}=1 + cup, \quad cup=2 + 3; \quad \textcolor{inc}{four}=?}
\end{center}
Only the label changes: \texttt{cup} is still 5 and the queried variable is still 6.
Answering 4 instead follows the name. We call 4 the \emph{name-suggested value} and that
answer a \emph{name error}. Because 4 occurs nowhere in the equations and differs from
every variable's value, this error can be distinguished from copying a digit.

Can asking the model to calculate before answering reduce this error? In our
\emph{answer-only} condition, the model is prompted to return just the final digit.
In our \emph{chain-of-thought} (CoT) condition, it is prompted to restate and evaluate
the equations before giving that digit. Here the calculation includes
\texttt{cup=2+3=5} and \texttt{four=1+5=6}. The comparison concerns what the model
generates; an answer-only response can still involve internal computation.

Misleading names already reduce accuracy on code tasks
\cite{le2025names,lam2025codecrash,le2026obfuscated}. We use arithmetic to make the
competing answers explicit, building on \citet{kudo2026faithful}'s study of computation
during CoT. We add a name that suggests a different answer while preserving the equations.
This conflict, like those in Stroop and shape--texture studies
\cite{stroop1935,geirhos2019texture}, separates the name's suggestion from the computed value.

Our first result is behavioral: CoT usually reduces name errors, although other accuracy
effects remain. We then ask whether the correct value becomes easier to recover inside
the model. A \emph{linear probe} is a small predictor trained on separate problems to
predict a variable's true digit from the model's internal activations. Recovering 6 with a
probe is a different observation from the language model answering 6. The probe--error
relationship depends on one low-accuracy model, so it does not explain the reduction.

\begin{figure*}[t]
\centering
\IfFileExists{figures/fig1_master_L3.pdf}{\includegraphics[width=\textwidth]{figures/fig1_master_L3.pdf}}{\fbox{\parbox{0.95\textwidth}{\centering\small\textcolor{red}{[Figure 1 not generated]}}}}
\caption{Eight panel models, queried-variable effects at level~3: \textbf{(a)} congruent-name accuracy gain and
\textbf{(b)} excess name errors, averaged over three demonstration sets with 95\% intervals.
Grey marks the two-point margin. \textbf{(c)} Maximum-over-layers true-value accuracy and
name-suggested prediction rate at each token position, pooled under a gold chain.
The value readout concerns the queried variable. Shading marks name positions; the dotted line
marks the position before its value.}
\label{fig:master}
\end{figure*}

\section{Study design}

We use \citeauthor{kudo2026faithful}'s single-digit arithmetic task: equations with addition
or subtraction, followed by a query. Our main level requires two operations, with the queried
variable defined before its dependency. Two models are also tested on all five levels
(Appendix~\ref{app:task}). Each problem has matched versions differing in one variable name
(Table~\ref{tab:conditions}). A congruent number word denotes that variable's true value;
an incongruent word denotes a different digit. We vary the queried and intermediate variables separately.
The name-suggested value cannot equal any variable's value or input digit. Names occupy one token in every
context, keeping twins' token positions identical.

\begin{table}[t]
\centering\small
\resizebox{\columnwidth}{!}{\begin{tabular}{@{}ll@{}}
\toprule
Name & Problem (answer: 6) \\
\midrule
Neutral & \texttt{pen=1 + cup, cup=2 + 3; pen=?} \\
Congruent & \texttt{\textcolor{con}{six}=1 + cup, cup=2 + 3; six=?} \\
Incongruent & \texttt{\textcolor{inc}{four}=1 + cup, cup=2 + 3; four=?} \\
\bottomrule
\end{tabular}}
\caption{Renaming the queried variable changes one cue while preserving the arithmetic.
The incongruent name denotes 4, which appears nowhere else in the problem.}
\label{tab:conditions}
\end{table}

Models either write a chain that restates and evaluates each equation or answer directly.
Every behavioral comparison uses three fixed sets of three worked examples. We evaluate
\NUM{n-models} checkpoints from the Llama, Gemma and OLMo families; the model list and tuning
comparisons are in Appendices~\ref{app:task} and~\ref{app:kind}.
The main behavioral table follows the preregistered 90\% mean neutral-CoT accuracy floor;
the excluded model's results remain in the appendix.

Our main measure is \emph{excess name errors}: the rate of answers matching the misleading
name's digit, minus the neutral twin's rate of answering that same digit. Subtracting this baseline
matters when accuracy is low. We also report congruent-minus-neutral accuracy
(\emph{facilitation}) and neutral-minus-incongruent accuracy (\emph{interference}). We call an effect reliable when its sign agrees across all three demonstration sets and its 95\%
cluster-bootstrap interval excludes zero. We call an effect small enough to meet the equivalence bound when its 90\% interval fits
wholly inside $\pm\NUM{eq-delta}$ points. We resample matched problems with their seed replicates together
(Appendix~\ref{app:equivalence}).

A linear probe predicts a variable's digit from a hidden state at one token and layer.
Probes train on separate neutral problems, with three training seeds. A control assigns a
random digit to each name, testing how well the same probe can learn name identity \cite{hewitt2019control}. Patching
replaces activations at the name's tokens with a matched twin's
\cite{zhang2024patching}. We report both a neutral-word control for disruption and an
alternative-name control for name-specific effects. Most internal measurements supply the
correct chain token by token. We also train on generated neutral chains, selecting layers on
computation-disjoint validation problems before testing held-out generations
(Appendix~\ref{app:round6-generated}).

\section{Results}

\begin{table*}[t]
\centering\small
\resizebox{\textwidth}{!}{\tabinput{behavior}}
\caption{Level~3, means over three demonstration sets (ranges in brackets). \emph{Helps}:
congruent-minus-neutral accuracy. \emph{Errors}: excess name errors. Both are percentage points.
$^{*}$ marks a consistent sign across sets and a 95\% interval excluding zero. CoT name-error effects
meet the two-point equivalence bound. The excluded low-accuracy model appears in
Table~\ref{tab:excluded}; additional models and levels appear in Tables~\ref{tab:kind}
and~\ref{tab:levels-appendix}.}
\label{tab:behavior}
\end{table*}

\subsection{Writing a chain limits name errors}

Table~\ref{tab:behavior} shows the naming effects at the main level. Without CoT, a congruent
name raises Llama-3.1-8B's accuracy by \NUM{facilitationllama318bdirectv1} points, and a
misleading name produces \NUM{lure_excessllama318bdirectv1} points of excess name errors on the queried
variable. Across levels, the largest effects occur where direct-answer accuracy is near chance; the matched baseline prevents those arbitrary digit answers from being attributed to
the name.

With CoT, \NUM{eq-cot-equivalent} of \NUM{eq-cot-cells} name-error contrasts meet the equivalence
bound. The exception is \NUM{eq-exc-model}: its intermediate-variable excess is
\NUM{eq-exc-mean} points [\NUM{eq-exc-lo}, \NUM{eq-exc-hi}]. Its neutral CoT accuracy is
\NUM{eq-exc-neutral-acc}\%, compared with \NUM{eq-exc-letter-acc}\% under letter names.
CoT therefore limits the name's numeric influence without removing every naming effect:
\NUM{cot-acc-cells} accuracy contrasts remain reliable
(Appendix~\ref{app:equivalence}).



\begin{figure}[t]
\centering
\IfFileExists{figures/fig3_patching_llama32-3b_L3_v1.pdf}{\includegraphics[width=\columnwidth]{figures/fig3_patching_llama32-3b_L3_v1.pdf}}{\fbox{\textcolor{red}{[Patching figure not generated]}}}
\caption{Llama-3.2-3B without CoT, queried variable: name errors removed by the neutral-twin patch, alongside
neutral-word damage and the alternative-name control, by layer.}
\label{fig:patch}
\end{figure}

\subsection{Value readouts and task competence}

A name's digit can be readable without representing the computed value. For Llama-3.2-3B,
the intermediate value is partly decodable after its defining equation (neutral accuracy
\NUM{p2-acc}; selectivity \NUM{p2-sel}). At the query, the name-identity control scores
\NUM{p3-ctl}. By the value-writing step, true-value accuracy reaches \NUM{p4-acc} or more
and name-suggested predictions fall to zero. A Gemma token sweep also finds different readout timing
for queried and intermediate values (Appendix~\ref{app:probe-extensions}).

Across \NUM{competence-models} models and \NUM{vs-cells} variable pairs, the \NUM{vs-decodable}
pairs with neutral probe accuracy of at least \NUM{vs-decodable-min} have excess name errors within
\NUM{vs-decodable-maxexcess} points of zero. The exception's intermediate readout is
\NUM{vs-undecodable-acc}. Yet the descriptive correlation between model-mean misleading-name
probe accuracy and excess name errors reverses from $r=\NUM{competence-probe-lure-r}$ to
$\NUM{competence-probe-lure-without-olmo-r}$ when OLMo-2-1B is omitted
(Appendix~\ref{app:round5-competence}). Overall task competence remains a competing account.
Within-problem fits also disagree on the association's sign (Appendix~\ref{app:sweeps}).

Gold-trained probes score \NUM{generated-olmo1-cal-lo}--\NUM{generated-olmo1-cal-hi}\% on
OLMo-2-1B's neutral generated chains. Training on generated neutral chains raises this to
\NUM{r6-generated-olmo1-cal-min}--\NUM{r6-generated-olmo1-cal-max}\%, still below the 90\%
calibration threshold in every cell. Llama-3.2-3B reaches
\NUM{r6-generated-llama3-cal-min}--\NUM{r6-generated-llama3-cal-max}\% with the same design,
but its name-error cases remain rare. The exception's generated-chain readouts therefore
remain descriptive (Appendix~\ref{app:round6-generated}).

Merely stating the correct value does not explain the pattern. We split the model's generated
chains by whether they state that value and compare only twins falling in the same split.
The exception still has \NUM{vw-olmo-excess} points of excess name errors
[\NUM{vw-olmo-lo}, \NUM{vw-olmo-hi}] across \NUM{vw-olmo-n} matched observations from
\NUM{vw-olmo-seeds} demonstration sets, despite stating the correct value. The full control is
in Appendix~\ref{app:controls}.

\subsection{Patching localizes the answer-only effect}

For Llama-3.2-3B without CoT, replacing name activations with the neutral twin's removes
\NUM{lure-removed-3b} of name errors when all layers are patched (Figure~\ref{fig:patch}).
A span-by-layer sweep places the strongest effects at the
defining equation in layers \NUM{grid-layers}, where removal reaches \NUM{grid-top-removed}
(Appendix~\ref{app:figs}). The neutral-word control also damages \NUM{control-damage} of correct
answers. The patch identifies a region sensitive to the name; disruption limits claims about the restored computation.

The Llama-3.1-8B controls do not distinguish a name-specific cause: neutral and alternative-name
patches remove errors similarly. Under a gold CoT, injecting the misleading name into a neutral
run of the exception model moves \NUM{exc-inject} of answers to the name-suggested value, while the word control
changes \NUM{exc-ctlword} at the same layer. These interventions test only the selected models and positions; they do not establish a common causal mechanism across models.



\section{Related work}

\citet{kudo2026faithful} locate intermediate-value readouts during arithmetic CoT; our competing
names distinguish a computed value from a lexical cue. Earlier Stroop studies test a single next
token \cite{wang2026stroop,hu2026conflict}. Faithfulness studies perturb prompts, chains or
parameters \cite{turpin2023unfaithful,lanham2023measuring,chen2025reasoning,tutek2025unlearning};
we combine a known conflicting cue with internal readouts. Positional copying and output priors
can also explain arithmetic errors \cite{liu2026readout,prabhakar2024deciphering}; our copy
control is reported in Appendix~\ref{app:sweeps}.
A companion code study \citep{anon2026code} examines value writes in Python traces; this
paper tests arithmetic name errors and their association with value-step readouts.

\section{Conclusion}

Writing a chain limits answers suggested by misleading names in almost every tested arithmetic contrast.
A correct value stated in the chain is insufficient to explain that reduction. The probe
association depends on the low-accuracy exception, whose generated-chain probes fail calibration
even after training on generated neutral chains.
The behavioral finding therefore remains stronger than the proposed internal explanation.

\section*{Limitations}

The task uses short synthetic arithmetic chains and English number words. It does not establish
how misleading identifiers affect real code. Identifier-style names were tested at one level
only (Appendix~\ref{app:controls}). Models are at most 12B parameters and receive raw completion
prompts; instruction-tuned models' chat templates and long-form reasoning modes are untested.
Demonstration choice changes accuracy substantially, which is why every behavioral claim must
hold across three sets.

Linear probes reveal decodable information, not necessarily the representation used to generate
an answer. Most internal measurements supply the correct chain token by token. Generated-chain
probes trained on generated neutral chains meet the calibration threshold for Llama-3.2-3B
but fail it for OLMo-2-1B (Appendix~\ref{app:round6-generated}). These probes cover two models
and select layers on independent validation computations. The nine-model competence comparison is descriptive;
variable roles are not independent models, and omitting OLMo reverses the readout--error
correlation. It does not separate value availability from overall task competence.
Patching changes mixed representations and can damage
correct answers, limiting the causal conclusions. Finally, the chain repeats names and equations;
the value-only format changes performance and cannot be treated as an equivalent intervention
(Appendix~\ref{app:simple}).

\bibliography{refs}

\appendix

\section{Task, models and replication}
\label{app:task}
A problem contains single-digit variables, addition or subtraction, and a query. The five
levels vary dependency order, depth and the presence of an unused equation
(Table~\ref{tab:levels}). Three same-level worked examples precede each test problem.
The evaluated models are \NUM{model-list}, plus \NUM{n-ladder} derived checkpoints
(Appendix~\ref{app:kind}).

\begin{table}[t]
\centering\small\setlength{\tabcolsep}{4pt}
\resizebox{\columnwidth}{!}{\begin{tabular}{@{}clccc@{}}
\toprule
Level & Problem (neutral names) & Steps & Held & Distr. \\
\midrule
1 & \texttt{pen=1 + cup, cup=2; pen=?} & 1 & 1 & 0 \\
2 & \texttt{pen=2 + 3, cup=1 + pen; cup=?} & 2 & 0 & 0 \\
3 & \texttt{pen=1 + cup, cup=2 + 3; pen=?} & 2 & 1 & 0 \\
4 & \texttt{pen=1 + cup, cup=2 + 3,} & 2 & 1 & 1 \\
  & \texttt{\ \ box=4 + 5; pen=?} & & & \\
5 & \texttt{pen=1 + cup, cup=2 + box,} & 3 & 2 & 0 \\
  & \texttt{\ \ box=1 + 2; pen=?} & & & \\
\bottomrule
\end{tabular}}
\caption{The five levels of \citet{kudo2026faithful}: additions or subtractions needed
(\emph{Steps}), variables unresolved when first read (\emph{Held}), and unused equations
(\emph{Distr.}). At level~1 the intermediate value is stated rather than computed.}
\label{tab:levels}
\end{table}

We reimplement \citet{kudo2026faithful}'s task and probe recipe rather than using their code.
The shared choices are single-digit outputs, same-level demonstrations, a 10-way linear
residual-stream classifier, full-batch SGD at $10^{-3}$ for 10,000 epochs, and their
span-by-layer-window patching design. For replication we train on letter names and report
$t^{*}$, $t^{*}_{\mathrm{eq}}$, $\mathrm{Acc}^{\prec\mathrm{CoT}}$ and
$\mathrm{Acc}^{\succ\mathrm{CoT}}$ at $\tau=0.9$.
The Llama-3.2-3B value-writing positions agree with their Table~3: equation
\NUM{replication-eq-v1} for the queried variable and \NUM{replication-eq-v2} for the intermediate.
Our additions are the conflicting names, matched baselines, demonstration-set claim rule and
equivalence bound.



On the letter condition, pre-chain probe maxima for Llama-3.2-3B are
\NUM{replication-pre-v1} and \NUM{replication-pre-v2}, compared with 17.8 and 33.2 reported by
\citet{kudo2026faithful}. The value-writing positions agree; the accuracy difference may reflect
training details. Llama-3.1-8B's intermediate value becomes decodable at equation
\NUM{replication8b-eq-v2}, but its queried value becomes decodable at equation
\NUM{replication8b-eq-v1}, earlier than Llama-3.2-3B's equation \NUM{replication-eq-v1}.
The larger model can therefore expose that value before writing its final evaluation.
Full heatmaps appear in Appendix~\ref{app:figs}.

\begin{table*}[t]
\centering\small
\tabinput{behavior_excluded}
\caption{The level-3 model below the preregistered 90\% mean neutral-CoT accuracy floor,
excluded from the main behavioral table. Columns and claim markers follow Table~\ref{tab:behavior}.}
\label{tab:excluded}
\end{table*}

\section{Equivalence checks}
\label{app:equivalence}
Following preregistration Amendment~3, equivalence requires a 90\% cluster-bootstrap
interval wholly inside the $\pm\NUM{eq-delta}$-point margin (two one-sided 5\% tests).
We resample matched sets with their demonstration-seed replicates together, using 4,000
bootstrap draws. Other behavioral intervals remain at 95\%.
For CoT accuracy, \NUM{eq-acc-equivalent} of \NUM{eq-acc-cells} facilitation and interference
contrasts meet the same equivalence bound. Of the \NUM{cot-acc-cells} reliable accuracy
contrasts, \NUM{eq-acc-claimable-equivalent} are also equivalent within the margin and
\NUM{eq-acc-claimable-outside} are not; a detectable effect can still be small.
The written-value control uses the same matched pairs for its excess estimate and interval;
the exception model's estimate pools all three demonstration sets and clusters on matched set.
\begin{table}[h]
\centering\small
\tabinput{equivalence}
\caption{Equivalence within the two-point margin uses 90\% intervals; reliability uses
95\% intervals excluding zero and the same sign across all three demonstration sets.
Each accuracy cell is one facilitation or interference contrast.}
\end{table}
\IfFileExists{tables/exception_patching.tex}{%
\begin{table*}[t]
\centering\small
\resizebox{\textwidth}{!}{\tabinput{exception_patching}}
\caption{OLMo-2-1B-Instruct, CoT, \NUM{exc-controls-n} matched pairs per target, patching name tokens at all
layers. Removed: percentage of baseline name-suggested predictions removed by the neutral-twin patch.
Word damage: percentage of correct neutral predictions damaged by the word control.
Alt. removed: baseline name-suggested predictions removed by the alternative-name patch. Alt. followed:
percentage following that source's name-suggested value. Rates are percentages;
-- means undefined. These readouts force the gold chain, so their baseline errors differ
from errors in the model's own generated chains.}
\end{table*}
}{}
\section{Tokenizer checks}
\label{app:tok}
We test candidate names in five tokenizer contexts and retain words that form one token
without merging with a neighbor. For Llama~3, all ten number words and 95 of 112 neutral
candidates pass. Gemma and OLMo use the same lists; runtime checks confirm equal token lengths
and positions within matched groups. Spaces around binary operators prevent merges such as
\texttt{1-two}. No instance in the reported runs is excluded by the layout check.

\section{Additional naming and written-value controls}
\label{app:controls}

\paragraph{Written values.} We split generated chains by whether they state the target's true
value. An observation enters the excess estimate only when its neutral twin falls in the same
split. The estimate and its 95\% interval use those same paired observations. Pooled over
eligible cells, the excess is \NUM{vw-excess-wrote} points in \NUM{vw-cells-wrote} cells where
the value is stated and \NUM{vw-excess-not} in \NUM{vw-cells-not} where it is not. The exception
persists in the stated-value split (main text), so this split does not explain its effect.

\paragraph{Identifier-style names.} At level~3, \NUM{ident-models} models receive names such as
\texttt{q4} instead of \texttt{four}, under three demonstration sets. Without CoT, queried-variable
excess name errors range from \NUM{ident-direct-lure-lo} to \NUM{ident-direct-lure-hi} points;
\NUM{ident-direct-lure-claim} of \NUM{ident-direct-cells} cells are reliable. For number words
on the same models, the range is \NUM{word-direct-lure-lo} to \NUM{word-direct-lure-hi}, with
\NUM{word-direct-lure-claim} of \NUM{word-direct-cells} reliable. The numeric name-error effect is
therefore stronger for number words in this comparison; congruent identifier names can still
improve accuracy.

\paragraph{Word-class effects.} With CoT at level~4, a number word on the queried variable costs
Llama-3.2-3B \NUM{wordclass-cost-L4} accuracy points whether the name agrees with the value or
not. Those errors restate the wrong equation. This is distinct from answering the digit the
name denotes, which is why the name-suggested value result cannot be read as removal of every naming effect.

\paragraph{Unused variables.} A number word on the level~4 distractor changes accuracy by at
most \NUM{irr-acc-delta-max} points. Its no-CoT excess name errors range from \NUM{irr-direct-lo} to
\NUM{irr-direct-hi}; \NUM{irr-direct-claimable} of \NUM{irr-direct-cells} cells are reliable.
With CoT, the range is \NUM{irr-cot-lo} to \NUM{irr-cot-hi}, with
\NUM{irr-cot-claimable} of \NUM{irr-cot-cells} reliable. At the answer position, probe probability of the name-suggested value
is \NUM{bound-mass} for a computed variable and \NUM{unbound-mass} for the distractor; the
latter's true value has probe accuracy \NUM{distractor-acc}.

\section{Full sweeps and selectivity}
\label{app:sweeps}

\begin{figure*}[h]
\centering
\IfFileExists{figures/fig2_tokens_llama32-3b_L3_cot_v2.pdf}{\includegraphics[width=\textwidth]{figures/fig2_tokens_llama32-3b_L3_cot_v2.pdf}}{\fbox{\parbox{0.95\textwidth}{\centering\small\textcolor{red}{[Figure 2, generated by scripts/58\_token\_figure.py]}}}}
\caption{In the layout of \citeauthor{kudo2026faithful}'s Figure~2. Per-token probes for the
incongruent condition,
Llama-3.2-3B, level~3, CoT. Top: how often the best-layer probe reads the true value (green) or
the name-suggested value (red). Bottom: true-value accuracy by layer. The displayed tokens are from one example;
the rates pool problems with different values and names. Dashed line: start of the supplied
correct chain; dotted line: position before the target value; arrow: the supplied value token.
Shading marks the name's tokens.}
\label{fig:heat}
\end{figure*}

\paragraph{Mixed-effects check.} We compare the paired bootstrap with a mixed logit model,
\texttt{name\_error \textasciitilde{} incongruent + C(seed)}, with a random intercept for
matched set. Each twin's outcome refers to the same name-suggested digit. Both procedures use the same
fixed subsample of \NUM{mix-cap} sets per cell. Of \NUM{mix-cells} fitted cells,
\NUM{mix-agree} agree on whether the interval excludes zero; \NUM{mix-unfitted} additional
cells have no outcome variation. The mixed fit supports an effect in \NUM{mix-gained} cells
where the bootstrap does not and withdraws support in \NUM{mix-lost} cells where it does
(\NUM{mix-lost-cells}). Subsampling also removes support from \NUM{mix-sub-lost-n} starred
Table~\ref{tab:behavior} cell under both procedures (\NUM{mix-sub-lost-cells}). The mixed
model uses variational Bayes, so its interval is an approximate posterior interval. We retain
the full-data bootstrap as the primary estimate on the reported percentage-point scale.

\paragraph{Instance-level association.} Within each model, level, regime, variable and
position, we regress name errors on the probe margin, $\log p_{\mathrm{true}}-\log p_{\mathrm{name}}$,
where $p_{\mathrm{name}}$ is the probe probability of the name-suggested value,
with standard errors clustered on matched set. The layer is selected by neutral accuracy.
Of \NUM{link-fits} fits, \NUM{link-sig} have $p<.05$: \NUM{link-neg} coefficients are negative
and \NUM{link-pos} positive. The direction is not consistent, and \NUM{link-pos-topmodel-n}
positive results come from the model supplying most fits. This leaves the instance-level
association unresolved. The model-level pattern is also sensitive to panel composition
(Appendix~\ref{app:round5-competence}).

\paragraph{Name-error share.} For comparison with shape--texture studies \cite{geirhos2019texture},
we also divide name errors by the total of name errors and correct answers, excluding congruent trials.
Across \NUM{lureindex-direct-n} level~3 model--variable pairs, the index ranges from
\NUM{lureindex-direct-lo} to \NUM{lureindex-direct-hi} without CoT and from
\NUM{lureindex-cot-lo} to \NUM{lureindex-cot-hi} with CoT. The main text uses paired excess
because the index does not subtract how often the neutral model answers that digit.

\paragraph{Positional-copy control.} Across \NUM{copy-n-err} generated CoT name errors over
models, levels and demonstration sets, the number immediately preceding the answer is the
name-suggested value in \NUM{copy-lure-frac} of cases and the target's true value in \NUM{copy-true-frac}.
A trailing name-suggested value is therefore absent from many errors. This check constrains the positional-copy
explanation of \citet{liu2026readout}, but does not identify what was copied in each case.

Probes report true-value accuracy, name-suggested prediction rate, probability of the name-suggested value and margin, with name-identity
control accuracy and selectivity \cite{hewitt2019control}. P1 is the name token, P2 the end of
its defining equation, P3 the query, P4 the position before its value is written, and P5 the
position before the final answer. P1 is a reference, not evidence for a computed value.
At P3 the control is near~1, so a readout of the name-suggested value there may reflect name identity.
Figure~\ref{fig:heat} shows the CoT token sweep; Appendix~\ref{app:figs} contains the other
conditions and direct regime. Letter-trained replication uses the metrics in
Appendix~\ref{app:task}.

\begin{figure}[h]
\centering
\IfFileExists{figures/fig5_kudo3_llama32-3b_L3_cot_v1_lureerr_i0.pdf}{\includegraphics[width=\columnwidth]{figures/fig5_kudo3_llama32-3b_L3_cot_v1_lureerr_i0.pdf}}{\fbox{\parbox{0.95\columnwidth}{\centering\small\textcolor{red}{[Figure 4, generated by scripts/64\_kudo\_fig3.py]}}}}
\caption{In the layout of \citeauthor{kudo2026faithful}'s Figure~3, for problems this model
answers with the name-suggested value. Each cell is the probe's top-1 prediction at that layer and token: green
the variable's true value, red the value the model answered, amber the value its name denotes,
grey anything else. The first panel is the probe's most common prediction on neutral problems,
which is what it emits where it reads nothing; cells carry information only where they differ
from it.}
\label{fig:kudo3}
\end{figure}

\begin{figure*}[h]
\centering
\IfFileExists{figures/fig_own_chain_llama32-3b_L3_cot_v1.pdf}{\includegraphics[width=\textwidth]{figures/fig_own_chain_llama32-3b_L3_cot_v1.pdf}}{\fbox{\parbox{0.95\textwidth}{\centering\small\textcolor{red}{[own-chain figure, generated by scripts/67\_own\_chain\_figure.py]}}}}
\caption{Probes trained at the value step, applied to the model's own generated chain rather
than to a force-decoded gold chain, for the \NUM{ownchain-n} level-3 incongruent problems that Llama-3.2-3B
answered wrongly under CoT (the first \NUM{ownchain-n} of \NUM{ownchain-total} wrong answers
in \NUM{n-test-sets} instances). Probes are
trained on neutral problems at the step that writes the variable's value and applied at every
generated token and every second layer; colours as in Figure~\ref{fig:kudo3}, boxed columns the
name's own tokens. This small sample contains no name errors and is illustrative rather than an estimate of their rate.
The fixed-sample transfer analysis and separate additional name-error cases appear in
Appendix~\ref{app:round5-generated}.}
\label{fig:ownchain}
\end{figure*}

\section{How the models were tuned}
\label{app:kind}

We compare six checkpoints on the Llama-3.1-8B architecture: pretrained, instruction-tuned,
single-task LoRA, multi-task LoRA, TIES-merged adapters and reasoning-tuned. Adapters are folded
into the weights before evaluation. Shared tokenizers and identical matched problems permit
paired comparisons (Table~\ref{tab:kind}). A pretrained/instruct Gemma pair supplies a second
family comparison.

\begin{table}[h]
\centering\small
\resizebox{\columnwidth}{!}{\tabinput{modelkind}}
\caption{Level 3, no CoT, name errors above the matched twin's rate, mean over three
demonstration sets. \emph{Acc.} is neutral accuracy. $^{*}$ as in Table~\ref{tab:behavior}.
The bottom block repeats the pretrained-to-instruct step in a second family, Gemma-3-4B.}
\label{tab:kind}
\end{table}

The pretrained and reasoning-tuned Llama checkpoints have reliable positive queried-variable
excess name errors without CoT (\NUM{kind-base-v1} and \NUM{kind-reas-v1} points). The three task-tuned
variants differ from the instruction-tuned checkpoint by at most \NUM{kind-maxgap} points
across \NUM{kind-tuned-cells} comparisons. These are observed differences, not tests of
checkpoint-to-checkpoint equivalence. Every checkpoint's CoT excess is at most
\NUM{kind-cot-max} points in magnitude.

Task accuracy complicates interpretation. The multi-task model's reliable negative excess
(\NUM{kind-ftm-v1}) accompanies lower neutral accuracy (\NUM{kind-ftm-acc}\%, versus
\NUM{kind-it-acc}\%). The reasoning-tuned checkpoint has the lowest accuracy
(\NUM{kind-reas-acc}\%). It also requires a system instruction to enable long-form reasoning;
our raw-completion prompts leave that mode off. The result concerns the checkpoint under this
prompt, not reasoning mode in deployment.

Only the instruct Gemma effect passes the claim rule (\NUM{kind-gi-v1}, versus
\NUM{kind-gb-v1} pretrained). The difference is untested and both estimates lie within the
margin, so the pair cannot establish where in training name interference arises.

\section{Probe recipe and patching scope}
\label{app:recipe}

We check two implementation choices inherited from \citet{kudo2026faithful}: the probe optimizer and the name occurrences included in a patch.

\paragraph{Recipe.} We refit the level-3 probes for Llama-3.2-3B with L-BFGS, and with SGD on
standardized inputs, against the reported recipe.

\begin{table}[h]
\centering\small
\resizebox{\columnwidth}{!}{\tabinput{recipe}}
\caption{Best-layer accuracy for the true value, neutral instances, under three probe recipes.}
\label{tab:recipe}
\end{table}

All three recipes give the same qualitative result: the value is decodable where the chain writes it
(\NUM{e17-value-lo}--\NUM{e17-value-hi}) and not at the end of its defining equation
(\NUM{e17-end-lo}--\NUM{e17-end-hi}). The query position moves, from \NUM{e17-sgd-query-v2}
under the reported recipe to \NUM{e17-std-query-v2} under standardization, within a range of
\NUM{e17-query-lo}--\NUM{e17-query-hi} over recipes and variables. We already report that
position as the presence of the name rather than a decoded value, because its control task
scores \NUM{p3-ctl}, so this variation does not change the interpretation of that position.

\paragraph{Scope.} Patching every occurrence of the name includes the copy the model writes in
its own output. In the direct regime, restricting the patch to the name's occurrences in the prompt leaves the
intermediate variable essentially unchanged, since its name never appears in the answer: removal is
\NUM{e18-v2-all-removed}\% at layers~\NUM{e18-v2-all-best} under both scopes. For the queried
variable the best windows remove similar shares (\NUM{e18-v1-prompt-removed}\% at
layers~\NUM{e18-v1-prompt-best} prompt-only, against \NUM{e18-v1-all-removed}\% at
layers~\NUM{e18-v1-all-best}), the prompt-only patch removes less in the later windows, and the
two disagree most in the first window, where the prompt-only patch removes
\NUM{e18-v1-prompt-w03}\% and damages \NUM{e18-v1-prompt-w03damage}\% of correct answers. That
is mechanical: the queried variable's name also sits at the position we read, so patching its
prompt copies alone leaves the model reading a name whose early representation no longer
matches. The localization we report does not depend on the choice.

\section{A chain that writes values without equations}
\label{app:simple}

Our chain restates each equation before substituting into it, as
\citeauthor{kudo2026faithful}'s does. Their Tables~5--6 also report a format that writes only the
sub-results, \texttt{cup=5, pen=6}. That format is the arithmetic counterpart of a code trace,
which writes a variable's value and never the expression that produced it, so we ran it on the
panel at level 3 with three demonstration sets, with two predictions registered in advance for both Llama models: that
the excess name errors at the intermediate's value step would exceed the \NUM{eq-delta}-point
margin, and that the value would still be decodable there (probe accuracy at least 0.85), making
any failure one of readout.

Both predictions fail as registered. Llama-3.1-8B writes the name-suggested value at the intermediate's value
step \NUM{simple-l8b-v2} points above its twin, inside the margin. Llama-3.2-3B is at
\NUM{simple-l3b-v2}, which passes our claim rule, but across the three demonstration sets that
is \NUM{simple-l3b-v2-seeds}, so one set carries most of it. More informative is what the probes say. At
the step that writes the value, the true value is decodable at \NUM{simpleprobe-l8b-simple-v1}
and \NUM{simpleprobe-l8b-simple-v2} for Llama-3.1-8B and at
\NUM{simpleprobe-l3b-simple-v1} and \NUM{simpleprobe-l3b-simple-v2} for Llama-3.2-3B, against
\NUM{simpleprobe-l8b-cot-v1} and \NUM{simpleprobe-l8b-cot-v2} for the same model under the full
chain. Neutral accuracy falls with it, to \NUM{simple-l8b-acc}\% and \NUM{simple-l3b-acc}\%.

Removing the equations lowers both accuracy and true-value probe readouts. It therefore fails to isolate a readout error while preserving those measures. This result does not support a shared explanation for value-only arithmetic chains and code traces.

\section{Additional levels, models and demonstration sets}
Table~\ref{tab:levels-appendix} repeats Table~\ref{tab:behavior} for the other four levels.
The pretrained Gemma-3-4B, the base half of the Gemma pair, is reported in
Table~\ref{tab:kind} rather than in Table~\ref{tab:behavior}.

\paragraph{Data, compute and software.} Each level has \NUM{n-test-sets} matched sets, each
rendered in every condition, and \NUM{n-probe-train} separately generated problems for
training probes; every SGD probe is trained with \NUM{probe-seeds} seeds. Every run here is
inference and the probes are linear; the three adapter checkpoints of Appendix~\ref{app:kind}
were trained in a separate project, and that compute is not counted below. The recorded allocation time totals at most \NUM{compute-gpuh} GPU-hours on NVIDIA A30, H100 and H200 cards, including idle time and shared workers that also ran companion-code experiments.
The models are from the Llama~3 \citep{grattafiori2024llama3}, Gemma~3 \citep{gemmateam2025gemma3}
and OLMo~2 \citep{olmo2024olmo2} families, with the reasoning-tuned checkpoint from
\citet{bercovich2025nemotron}, each used under its release licence for research. Models run in
PyTorch \citep{paszke2019pytorch} with Transformers \citep{wolf2020transformers}, and the probes are
trained in PyTorch; the L-BFGS refit of Appendix~\ref{app:recipe} uses scikit-learn
\citep{pedregosa2011sklearn}, and the mixed-effects check uses statsmodels
\citep{seabold2010statsmodels}.

\begin{table}[h]
\centering\small
\resizebox{\columnwidth}{!}{\tabinput{behavior_levels}}
\caption{As Table~\ref{tab:behavior}, for levels 1, 2, 4 and 5 (Llama-3.2-3B and Llama-3.1-8B). Level~1 is the no-computation
control: the intermediate variable's value is stated rather than computed.}
\label{tab:levels-appendix}
\end{table}

Per-set values, the pooled estimate and the claim rule are released with the code, together
with the level-4 irrelevant-name-suggested value
control, the name-suggested value-distance analysis, the teacher-forced name-suggested prediction rates, and the error taxonomy
behind the word-class effect.

% Round-four probe extensions
\section{Additional value readouts}
\label{app:probe-extensions}

We extend the value-step comparison to Llama-3.2-3B-Instruct, Llama-3.1-8B-Instruct and
Gemma-3-12B-Instruct. Each probe trains on 10,000 separate neutral problems with three probe
seeds, using demonstration set 7. Table~\ref{tab:probe-extensions} reports the added pairs;
the main comparison includes them. We select layers by neutral selectivity, as in the original
analysis. The misleading-name accuracy uses the same selected layer. This extension tests
the aggregate association; it does not resolve the within-problem link.

\begin{table*}[t]
\centering\small
\tabinput{probe_extensions}
\caption{Additional probes at each variable's value-writing position. Neutral and misleading
columns are true-digit accuracy; control is neutral name-identity accuracy. Rates are percentages,
averaged over three probe seeds. The layer maximizes neutral accuracy minus control accuracy.}
\label{tab:probe-extensions}
\end{table*}

We also probe every token for Gemma-3-4B-Instruct (Figures~\ref{fig:gemma-query}
and~\ref{fig:gemma-intermediate}), using 4,000 independent neutral training problems and
1,000 test problems per condition. The tokens shown come from one example; the curves and
heatmaps pool problems with different names and values. Maximum-over-layers curves describe
where a digit can be read; they do not identify the representation used to answer.
On neutral problems, the best pre-chain accuracy is \NUM{gemma-token-pre-v1}\% for the
queried digit and \NUM{gemma-token-pre-v2}\% for the intermediate digit. The queried readout
\NUM{gemma-token-timing-v1}; the intermediate readout \NUM{gemma-token-timing-v2}.
Token zero is the first token of the chain; negative positions are in the prompt.

\begin{figure*}[t]
\centering
\includegraphics[width=\textwidth]{figures/fig2_tokens_gemma3-4b-it_L3_cot_v1.pdf}
\caption{Gemma-3-4B-Instruct, queried variable, misleading-name problems under the correct
chain. Top: maximum-over-layers true-value accuracy and name-suggested prediction rate; bottom: true-value accuracy
by layer. Shading marks the variable's name; the dotted line marks the position before its
value; the arrow marks the supplied digit itself.}
\label{fig:gemma-query}
\end{figure*}

\begin{figure*}[t]
\centering
\includegraphics[width=\textwidth]{figures/fig2_tokens_gemma3-4b-it_L3_cot_v2.pdf}
\caption{The same Gemma token sweep for the intermediate variable. Both figures use
neutral-trained probes; no misleading-name problems enter training.}
\label{fig:gemma-intermediate}
\end{figure*}

\input{tables/figure_dump}

\input{tables/round5_followups}
\input{tables/round6_followups}

\end{document}
